\section{Task 3 -- RQ3: Methodology}
\label{sec:rq3-methodology}

\textbf{RQ3:} To what extent does PlanPlay-SQL improve Qwen3-0.6B?

\subsection{Model and Option Selection}
\label{sec:rq3-model-selection}

\textbf{We select Option B -- Adaptation / Learning to Adapt: we adapt Qwen3-0.6B's parameters using PlanPlay-SQL}, rather than Option A's reasoning-only strategies. Text-to-SQL requires substantial schema understanding, and precise SQl generaiton. This makes Qwen3-0.6B an interesting model to improve, as its relatively small size makes it computationally efficient and accessible for fine-tuning.

\subsection{Approach: PlanPlay-SQL in Four Stages}
\label{sec:rq3-approach}

We call our method \textbf{PlanPlay-SQL}: the model first writes a short \emph{plan}, then is trained by verified self-\emph{play}. It combines three ideas from recent work on small-model Text-to-SQL into four stages, each motivated by a specific gap in the base model:

\begin{enumerate}
\item \textbf{Plan prompting.} From Enhancing LLM Fine-Tuning for Text-to-SQLs (2024) and FINER-SQL (2026), the model writes a short plan (tables and joins, columns, filters, aggregation and ordering) before the SQL. \emph{Why:} a model with no thinking mode otherwise has no scratch space to organize the schema before committing to SQL syntax.
\item \textbf{Supervised warm start}: we attach a fresh LoRA adapter to the frozen base checkpoint and train only the adapter to imitate the gold SQL and its rule-based plan. \emph{Why:} Task 1 showed many of the model's outputs are valid SQL that differs from the gold query, so the fastest way to close that gap is to show it Spider's own conventions directly.
\item \textbf{Verification-based iterative fine-tuning (VBI-FT)}, following SPFT-SQL (2025): the model answers training questions, we execute its SQL and the gold SQL, keep the answers with matching results as positives, and fine-tune on them. \emph{Why:} this trains on the model's own output distribution rather than only gold text, which should reduce remaining execution errors without needing additional human-labeled data.
\item \textbf{Self-play fine-tuning}, using the logistic loss of SPIN (Chen et al., 2024): we train the best checkpoint to prefer the gold answer over an incorrect answer produced by the worst checkpoint. \emph{Why:} SPFT-SQL (2025) found this stage added a further 0.6 to 0.8 points on top of verification-based fine-tuning, on much larger models, so we test whether a comparable gain appears at 0.6B.
\end{enumerate}

We train using Spider's provided training split, and checkpoints are selected on held-out training databases (\ref{sec:rq3-evaluation}), so the dev set stays untouched and evaluation stays cross-domain. We skip SPFT-SQL's (2025) synthetic question generation because Spider already provides gold SQL to execute against.

We expect PlanPlay-SQL to improve performance in several ways: stage 1 (the plan prompt with no fine-tuning) serves as a control that isolates what the prompt format contributes on its own, so any further gain from stage 2 onward can be attributed to training the model to use the format rather than to the format itself. Stage 2 (the warm start) should raise exact and execution accuracy most, since it directly targets the style gap identified in \ref{sec:rq3-model-selection}. Stage 3 (verification-based fine-tuning) then trains on the model's own execution-verified outputs rather than only gold text, which should reduce its remaining execution errors without needing additional gold-labeled data. Stage 4 (self-play) should add a further small gain on top of stage 3, by analogy with SPFT-SQL's (2025) 0.6-to-0.8-point gain on larger models. To measure each part, we report five rows on the same dev set and metrics as \hyperref[sec:rq1-results]{Task 1}: the original solution, the plan prompt only, the warm start only, warm start plus verification fine-tuning, and the full PlanPlay-SQL.



\subsection{Model Training}
\label{sec:rq3-training}

Spider's Gold SQL has no plan annotation, so we use a rule-based parser to generate stage 2's warm-start target. The parser reads each gold query's top-level clauses (FROM/JOIN, SELECT, WHERE, GROUP BY/ORDER BY/LIMIT, and any set operator) with a small regex-based SQL splitter and deterministically renders them into the five-line plan format used as the training target.

All stages use LoRA adapters (rank 32, alpha 64) on top of the frozen base checkpoint, rather than full fine-tuning, as LoRA fine-tuning achieves identical performance as full fine-tuning for most tasks at significantly less compute, memory, and storage cost.

Table~\ref{tab:rq3-training-settings} gives the learning rate, epochs, batch size, and training-item counts actually used for each stage; all stages share seed 0 and the same 707-question, 20-database held-out validation set described in \ref{sec:rq3-evaluation}. ``Samples per question'' (k) is the number of candidate SQL completions per training question sampled at temperature 0.7 and top-p 0.95 during verification-based and self-play fine-tuning (stages 3-4), before the correctness filter is applied; the resulting training-item count after filtering is reported in the ``training items'' column, since not every sampled question yields a usable positive or preference pair. Self-play additionally requires a second, frozen \emph{opponent} checkpoint: the opponent is sampled to supply the incorrect half of each preference pair, and the loss measures the trained (``policy'') checkpoint's margin between the correct and incorrect completions relative to the opponent's own margin on the same pair. Rows A-C (plain supervised or verification-based fine-tuning) have no opponent, since they have no preference pair at all. Row D is labeled as step 1 of the lower-LR self-play branch, but the step itself is a plain VBI-FT pass on a separate 800-question sample, used only to produce a slightly improved policy checkpoint for row E; it likewise has no opponent. Rows E and F are the only two rows that actually run the self-play preference-loss step.

\begin{table}[H]
\centering
\footnotesize
\caption{Training settings for each stage. Learning rate is the peak LR under a linear schedule. ``---'' marks columns that do not apply to a given row (no opponent checkpoint for non-self-play stages; no sampling for gold-only stages).}
\label{tab:rq3-training-settings}
\resizebox{\textwidth}{!}{%
\begin{tabular}{llllccccc}
\toprule
Stage & Policy init & Opponent & Training items & \# of Samples per Training Question & LR & Epochs & Batch & Rounds \\
\midrule
A: Warm start (plan format, 1500 gold) & Qwen3-0.6B & --- & 1500 gold SQL+plan pairs & --- & 2e-4 & 3 & 32 & 1 \\
B: A + 1500 further gold (supervised control) & Warm start & --- & 1500 gold SQL+plan pairs & --- & 1e-4 & 1 & 32 & 1 \\
C: B + VBI-FT (1500 q, hard only) & Warm start & --- & 476 verified positives (of 1500 sampled) & 4 & 1e-4 & 1 & 32 & 1 \\
D: self-play (lower LR), step 1/2: internal VBI-FT (own 800 q) & Warm start & --- & 946 verified positives (of 800 sampled) & 4 & 5e-5 & 1 & 32 & 1 of 2 planned \\
E: self-play (lower LR), step 2/2: preference training & D's checkpoint (val.\ acc.\ 0.624) & Warm start & 279 verified preference pairs (of 800 sampled) & 4 & 1e-5 & 1 & 32 & 1 of 2 planned \\
F: self-play (higher LR), round 1 & Warm start & Warm start (same checkpoint; round 1 has no earlier opponent) & 727 verified preference pairs (of 1500 sampled) & 4 & 1e-4 & 2 & 32 & 1 \\
\bottomrule
\end{tabular}}
\end{table}

\subsection{Evaluation Setup}
\label{sec:rq3-evaluation}

See \hyperref[sec:appendix4]{Appendix 4} for the full PlanPlay-SQL prompt template and the rule-based parser that renders gold SQL into the plan format used as the warm-start training target.

We evaluate on the same 1034-question Spider dev set and the same metrics as \hyperref[sec:rq1-results]{Task 1} (component matching, exact matching, execution accuracy), so that Task 3's comparison is a direct extension of Task 1 rather than a new benchmark. The adapted model uses the PlanPlay-SQL prompt, greedy decoding, and a maximum of 384 new tokens (Task 1: 256).

Checkpoints across all stages are selected on a held-out validation set of 707 questions from 20 training databases that were never used for training, keeping the dev set untouched for a single final comparison as well as keeping evaluation cross-domain. This mimics how the Spider benchmark was used in Task 1.

\section{RQ3: Results}
\label{sec:rq3-results}

\subsection{Overall Results}
\label{sec:rq3-results-overall}

Table~\ref{tab:rq3-dev-results-06b} compares the original Qwen3-0.6B from \hyperref[sec:rq1-results]{Task 1} with the PlanPlay-SQL-adapted model on the same 1034 Spider dev questions, with the same metrics and the same execution-checking databases.

\begin{table}[H]
\centering
\footnotesize
\caption{Original and adapted Qwen3-0.6B on the Spider dev set, overall and by difficulty. Component matching is accuracy / recall / F1 (macro-averaged as in \hyperref[sec:rq1-interpretation]{Task 1}). Execution accuracy is on mock databases.}
\label{tab:rq3-dev-results-06b}
\resizebox{\textwidth}{!}{%
\begin{tabular}{llccccc}
\toprule
Metric & Model & easy & medium & hard & extra & all \\
\midrule
Problems (n) & & 248 & 446 & 174 & 166 & 1034 \\
Component matching & Original (Task 1) & 0.675 / 0.308 / 0.423 & 0.582 / 0.235 / 0.334 & 0.540 / 0.141 / 0.223 & 0.578 / 0.160 / 0.250 & 0.628 / 0.209 / 0.314 \\
 & Adapted & 0.809 / 0.763 / 0.786 & 0.750 / 0.643 / 0.693 & 0.815 / 0.626 / 0.708 & 0.696 / 0.478 / 0.567 & 0.783 / 0.636 / 0.702 \\
Exact matching & Original (Task 1) & 0.157 & 0.137 & 0.006 & 0.006 & 0.099 \\
 & Adapted & 0.762 & 0.594 & 0.385 & 0.265 & 0.546 \\
Execution (mock DB) & Original (Task 1) & 0.173 & 0.179 & 0.029 & 0.036 & 0.130 \\
 & Adapted & 0.794 & 0.583 & 0.443 & 0.277 & 0.561 \\
\bottomrule
\end{tabular}}
\end{table}

The adapted model is better than the original on every metric and every difficulty level. Exact matching accuracy rises from 0.099 to 0.546 overall, execution accuracy on the mock databases rises from 0.130 to 0.561, and component matching recall rises from 0.209 to 0.636, so the model now writes most of the components the gold query contains. Our motivating claim in \ref{sec:rq3-model-selection} (adapting Qwen3-0.6B's parameters (Option B) on the Spider training split would close its style gap) held, with much larger gains than we expected from a 0.6B model.

\subsection{Results by Difficulty Level}
\label{sec:rq3-results-difficulty}

The gain is not uniform across difficulty. The largest relative gains are on hard (exact matching 0.006 to 0.385) and extra (0.006 to 0.265) problems, which the original model almost never solved. This is consistent with the warm start fixing a systematic style/convention gap rather than a difficulty-specific one, since harder questions have the most schema and clause structure for the plan format to organize. Easy and medium problems also improve substantially (exact matching 0.157 to 0.762 and 0.137 to 0.594 respectively), so the adaptation is not merely redistributing accuracy toward harder cases at the expense of easy ones.

\subsection{Metrics Breakdown}
\label{sec:rq3-metrics-breakdown}

The three metrics agree on direction but differ in what they reveal. Execution accuracy (0.130 to 0.561) and exact matching (0.099 to 0.546) rise by nearly the same amount, so the model is not merely producing lucky executions on the mock databases but genuinely matching gold queries. Component matching recall rises further in relative terms (0.209 to 0.636), meaning the model now includes most gold-query components even in cases where it falls short of an exact or executable match (see Table~\ref{tab:rq3-component-breakdown-06b})

\begin{table}[H]
\centering
\footnotesize
\caption{Component Matching accuracy by SQL component, original vs.\ PlanPlay-SQL-adapted Qwen3-0.6B, on the same 1034 Spider dev questions as table~\ref{tab:rq3-dev-results-06b}. Component-matching accuracy does not depend on the execution databases.}
\label{tab:rq3-component-breakdown-06b}
\resizebox{0.7\textwidth}{!}{%
\begin{tabular}{llccccc}
\toprule
Component & Model & easy & medium & hard & extra & all \\
\midrule
select & Original & 0.844 & 0.819 & 0.960 & 0.862 & 0.843 \\
 & Adapted & 0.956 & 0.855 & 0.975 & 0.792 & 0.892 \\
select (no AGG) & Original & 0.875 & 0.839 & 0.960 & 0.862 & 0.861 \\
 & Adapted & 0.973 & 0.866 & 0.975 & 0.792 & 0.901 \\
where & Original & 0.711 & 0.359 & 0.150 & 0.158 & 0.391 \\
 & Adapted & 0.875 & 0.833 & 0.694 & 0.569 & 0.784 \\
where (no OP) & Original & 0.711 & 0.391 & 0.200 & 0.368 & 0.438 \\
 & Adapted & 0.896 & 0.833 & 0.726 & 0.667 & 0.808 \\
group (no Having) & Original & 1.000 & 0.667 & 1.000 & 1.000 & 0.900 \\
 & Adapted & 0.850 & 0.741 & 0.857 & 0.696 & 0.758 \\
group & Original & 0.000 & 0.667 & 1.000 & 1.000 & 0.800 \\
 & Adapted & 0.750 & 0.672 & 0.857 & 0.625 & 0.696 \\
order & Original & 0.810 & 0.571 & 0.000 & 0.333 & 0.552 \\
 & Adapted & 0.864 & 0.833 & 0.900 & 0.875 & 0.866 \\
and/or & Original & 1.000 & 0.916 & 0.893 & 0.883 & 0.927 \\
 & Adapted & 1.000 & 0.966 & 0.959 & 0.905 & 0.964 \\
IUEN & Original & 0.000 & 0.000 & 0.000 & 0.000 & 0.000 \\
 & Adapted & 0.000 & 0.000 & 0.400 & 0.231 & 0.288 \\
keywords & Original & 0.800 & 0.591 & 0.240 & 0.310 & 0.572 \\
 & Adapted & 0.928 & 0.906 & 0.803 & 0.811 & 0.878 \\
\bottomrule
\end{tabular}}
\end{table}

\subsection{Stage-wise Ablation}
\label{sec:rq3-ablation}

To see which part of the method helps, table~\ref{tab:rq3-validation-accuracy-06b} reports accuracy on the held-out validation set of 707 questions after each stage. The rows correspond to table~\ref{tab:rq3-training-settings} as follows (counting table~\ref{tab:rq3-validation-accuracy-06b}'s first data row, ``Base model, original prompt*,'' as row 1): the warm start row 3, the supervised control is row 4, VBI-FT is row 5, the combined lower-LR self-play branch is row 6, and the higher-LR self-play branch is row 7.

\begin{table}[H]
\centering
\footnotesize
\caption{Validation accuracy (execution match on mock databases, 707 questions) after each stage. The row marked * uses a 300-question subset of the same validation set.}
\label{tab:rq3-validation-accuracy-06b}
\begin{tabular}{lc}
\toprule
Stage & Validation accuracy \\
\midrule
1. Base model, original prompt* & 0.407 \\
2. Base model, plan prompt, no training* & 0.050 \\
3. Warm start on 1500 gold questions (plan format) & 0.622 \\
4. + 1500 further gold questions (supervised control) & 0.652 \\
5. + VBI-FT on the same 1500 questions (no gold SQL, hard questions only, 476 verified samples) & 0.651 \\
6. + self-play, lower learning rate (1e-5, 9 steps) & 0.631 \\
7. + self-play, higher learning rate (1e-4, 46 steps) & 0.034 \\
\bottomrule
\end{tabular}
\end{table}

Stage 1's premise did not hold: giving the base model the plan-prompt format with no fine-tuning does not add scratch space for it to organize the schema, as \ref{sec:rq3-approach} motivated it. Instead, it collapses validation accuracy from 0.407 to 0.050, a far larger effect than a neutral control would produce. A model never trained on the plan format has no reason to follow it, so it most likely spends the budget on a malformed or incomplete plan instead of SQL, rather than using the extra structure productively. 

However, the warm start recovers from this collapse and accounts for most of the overall gain (0.050 to 0.622, or 0.407 to 0.622 against the better zero-shot baseline). Verification-based fine-tuning (stage 3) matched the supervised control (0.651 vs.\ 0.652, a difference far smaller than the roughly 2-point sampling error of 707 questions) while using only 476 self-generated, execution-verified training items instead of 1500 gold ones and no gold SQL text. This shows the model can close most of its remaining gap on its own verified outputs using LoRA fine-tuning and VBI-FT.


Our expectation that stage 4 would add a further, smaller gain as in SPFT-SQL (2025) did not hold. Self-play did not help: with a low learning rate it barely changed the model (0.622 to 0.631), and with a higher learning rate the training loss fell from 0.69 to 0.25 but accuracy collapsed to 0.034 with most outputs failing to execute. This is more pronounced than SPFT-SQL's (2025) own small gains on much larger models: self-play's margin at 0.6B is thin enough that a higher learning rate overshoots it entirely. We did not rerun the higher-learning-rate setting with intermediate checkpointing or a lower learning rate, so we cannot rule out that a different schedule would recover some of the gain the lower-LR run showed.

\subsection{Examples}
\label{sec:rq3-examples}

Out of 1034 dev questions, the adapted model gets 474 right that the original got wrong, loses 28 that the original got right, and both get 106 right (execution match on mock databases). One fixed example is the network\_1 question ``Show the student IDs and numbers of friends corresponding to each.'' The original model wrote \texttt{SELECT h.ID, h.name, h.grade, f.friend\_id FROM Highschooler h JOIN Friend f ON h.ID = f.student\_id}, which lists friends instead of counting them. The adapted model first wrote the plan \texttt{tables: friend; select: student\_id , count(*); group/order: group by student\_id} and then \texttt{SELECT student\_id , count(*) FROM friend GROUP BY student\_id}, which matches the gold query. One broken example is the tvshow question ``What are the titles of the cartoons sorted alphabetically?'' The original model wrote a correct query (\texttt{SELECT Title FROM Cartoon ORDER BY Title ASC}), but the adapted model's plan misspelled the table name (\texttt{tables: CARTON}), and the SQL that followed inherited the error (\texttt{SELECT Title FROM CARTON ORDER BY Title}), which fails to run. In this case the plan step propagated a mistake into the final query.

\subsection{Limitations}
\label{sec:rq3-limitations}

\begin{itemize}
\item \textbf{Mock databases.} The official Spider database files with rows were not available, so we generated databases with random rows seeded from the literals in each question's gold query. Execution accuracy on them only approximates the official execution accuracy. For example, Task 1's execution accuracy for the original model drops from 0.209 on schema-only databases to 0.130 on the mock databases (table~\ref{tab:overall-performance}), because on empty databases every query that returns nothing counts as "correct". Exact matching accuracy does not depend on the databases.
\item \textbf{Different prompt and token budget.} The adapted model uses a different prompt and a larger token budget than the original, so the gain combines fine-tuning, the plan format, and the longer limit.
\end{itemize}
